\section{Experimental Design}
\label{sec:experimental-methodology}

\subsection{Objectives and analyses}
\label{sec:rqs}

Objective~1 is met by the harness in \Cref{sec:method}; the remaining three map
onto the analyses as follows. Objective~2 has two parts---the level at each size
and the scale interaction---which have equal standing and are tested by two separate analyses
(per-size and difference-in-differences, \Cref{sec:stats}); significance at one
size must not be read as evidence about downscaling. Objective~3 applies the same
two analyses to the alignment panel and stops at measurable signals; construct
validation requires expert evidence beyond automated scores. Objective~4 is tested
by the three ablation arms below, which run inside the same ladder sessions.

\subsection{Comparison conditions}
\label{sec:conditions}

The primary comparison includes seven systems: AlignmentGraph-ISD and six
benchmark agents (Baseline, EduPlanner, ADDIE-Agent, RPISD-Agent,
Dick-Carey-Agent, and ReAct-ISD) at every size and every run.

For the ablation study, three arms of the proposed system run alongside, each
disabling exactly one part of the mechanism in \Cref{sec:method}: $-$Verifier
disables the check-and-repair layer (\Cref{sec:qc}) while keeping decomposition
and graph context; $-$GraphContext disables the context-slice reads $C(s)$
while keeping writes under the contract $W(s)$ (\Cref{sec:alignment-graph});
and the skeleton disables both, keeping only specialist decomposition. This
design tests the two mechanisms built on top of decomposition, not
decomposition itself---the inherited architectural foundation
(\Cref{sec:decomposition}). All three arms run inside the same ladder sessions,
so arm-minus-A0 deltas share a judge session and avoid judge drift; the
in-session A0 also serves as a consistency check against the three primary A0
runs.

\subsection{Data and model ladder}
\label{sec:dataset}
\label{sec:setup}

Every condition runs on the \texttt{test\_90} split of
ISD-Agent-Bench~\cite{isdagentbench}: 90 scenarios selected by stratified
sampling with seed 42, shared by all systems. The ladder comprises Qwen
checkpoints at 0.8B, 2B, 4B, and 9B; holding the model family constant means
only capacity changes between sizes. Each agent--size cell has three
independent runs over all 90 scenarios. Scenarios and prompts for all seven
agents are English by design; the only baseline-specific intervention
translates ReAct-ISD's prompt scaffolding into English without changing its
method. Manifests record checkpoints, prompts, seeds, decoding settings, and
judges.

\subsection{Deployment envelope and serving conditions}
\label{sec:assumptions}

The deployment target is a single 12--16\,GB GPU. At two bytes per parameter,
bf16 weights require approximately 1.6\,GB for 0.8B, 4\,GB for 2B, 8\,GB for
4B, and 18\,GB for 9B; the 9B model therefore does not fit, and 4B is the
largest in-target checkpoint. The remaining VRAM must hold the KV cache and
runtime overhead, leaving a budget of roughly 8k tokens per call. This chain of
reasoning is checkable independently of experiment timing; 9B is an
out-of-target reference point, not a proposed deployment configuration.

All agents run through vLLM with streaming enabled, thinking disabled, and a
common generation cap. We cap generation at 8192 tokens per call---the budget a
4B model leaves on a single 12--16\,GB GPU, the deployment target this work
claims. The 24--32\,GB pods were rental convenience for running the ladder in
parallel, not the claimed envelope. Whether an architecture completes within
this limit is therefore a measured deployment property.

\subsection{Measures}
\label{sec:metrics}

Two measurement families are applied to the same saved outputs: the benchmark's
LLM-judge rubric scores, and a deterministic alignment-signal measurement that
uses no LLM. They are mechanically independent, so agreement between them carries
more weight than either family alone.

\subsubsection{Judge-based measures}

For Objective~2, two fixed judges from families separate from the agents---
\texttt{openai/gpt-4o-mini} and \texttt{google/gemini-2.5-flash-lite}---score
every artifact; this separation reduces self-preference
risk~\cite{panickssery2024llm}. The scenario score is the mean of the two
judges. The three judge signals are reported in parallel with identical
analyses and equal standing (\Cref{tab:judge-signals}); ADDIE leads the
narrative because its rubric matches the construct Objective~2 addresses.

\begin{table}[t]
  \caption{The three judge signals (scale /100).}
  \label{tab:judge-signals}
  \small
  \begin{tabular}{@{}lp{0.35\columnwidth}p{0.24\columnwidth}@{}}
    \toprule
    Signal & Measured construct & Role \\
    \midrule
    \texttt{addie\_median} & content quality of the design (artifact rubric) &
    leads the narrative \\
    \texttt{trajectory\_score} & tool-use process conformance (BFCL) &
    reported in parallel \\
    \texttt{total\_score} & the benchmark's $0.7\,$ADDIE$+0.3\,$Traj composite &
    comparability with the benchmark \\
    \bottomrule
  \end{tabular}
\end{table}

\subsubsection{Alignment-signal measures}
\label{sec:alignment-measures}

The judge scores above are produced by LLMs, whose rubric-scoring biases are
systematically documented~\cite{zheng2023judging,panickssery2024llm,norman2026reliability},
and here the judge and generator families overlap. The study therefore adds a
second, mechanically independent measurement that reads the semantic layer
corresponding to the structural relations the graph makes explicit: it measures
signals associated with the three facets of \Cref{eq:aligned} on the rendered
output of every system. It operationalizes established criteria and follows two
principles:
\begin{itemize}
  \item \emph{No LLM in measurement}: there is no prompting, sampling, or matching
  threshold; after the encoder and Bloom classifier are pinned, every score is a
  fixed function of text.
  \item \emph{Read only text from saved output}: self-declared graph identifiers
  are excluded from score computation and retained only as validity evidence, so
  every system is measured on the same basis.
\end{itemize}

The panel contains seven signals from two mechanically independent instrument
families: \emph{Family A} measures textual correspondence between artifacts with
embeddings; \emph{Family B} measures cognitive demand through Bloom levels and
does not use cosine similarity. From each system output, the protocol extracts four
text sets: $\mathcal{O}$ (objectives), $\mathcal{A}$ (assessment items),
$\mathcal{X}$ (learning activities), and $\mathcal{E}$ (evaluation items). Family
A uses the mean--max rectified-cosine operator
\[
  \mathrm{sim}(S,T)
  = |S|^{-1}\sum_{u\in S}\;\max_{t\in T}\;\max\bigl(0,\cos(e_u,e_t)\bigr),
\]
where $e_u$ is the embedding of text span $u$; $\max(0,\cdot)$ is a scale choice
on $[0,1]$, not a matching threshold. \Cref{tab:panel-signals} specifies the seven
panel signals: four Family A signals correspond to the three facets of
\Cref{eq:aligned} plus a reverse diagnostic direction, while three Family B signals
directly inherit established alignment instruments from educational measurement.

\begin{table*}[t]
  \caption{The seven panel signals. Family A applies the $\mathrm{sim}$ operator
  to four text sets extracted from output; Family B uses Bloom levels and does not
  use cosine similarity. \dag{} = non-directional and excluded from directional
  counts under the panel interpretation rule. Parenthetical names are the
  corresponding fields in the generated artifact.}
  \label{tab:panel-signals}
  \small
  \begin{tabular}{@{}lp{0.20\textwidth}p{0.22\textwidth}p{0.27\textwidth}@{}}
    \toprule
    Symbol & Formula & Measurement & Selection basis \\
    \midrule
    $\sigma_{\mathrm{meas}}$ {\scriptsize(objective\_assessment\_similarity)} &
    $\mathrm{sim}(\mathcal{O},\mathcal{A})$ &
    how closely an assessment corresponds in content to each objective &
    \emph{measures} facet of \Cref{eq:aligned}; precedent for embedding-based
    item--standard mapping~\cite{xu2025itemalignment,clelandhuang2018traceability} \\
    $\sigma_{\mathrm{supp}}$ {\scriptsize(objective\_activity\_similarity)} &
    $\mathrm{sim}(\mathcal{O},\mathcal{X})$ &
    how closely learning activities address each objective &
    \emph{supports} facet; constructive alignment requires activities to support
    objectives~\cite{biggs1996enhancing} \\
    $\sigma_{\mathrm{eval}}$ {\scriptsize(objective\_evaluation\_similarity)} &
    $\mathrm{sim}(\mathcal{O},\mathcal{E})$ &
    how closely evaluation covers each objective &
    \emph{evaluates} facet; Quality Matters requires evaluation to connect to
    objectives~\cite{qm2018rubric} \\
    $\sigma_{\mathrm{rev}}$\,\dag{} {\scriptsize(assessment\_objective\_similarity)} &
    $\mathrm{sim}(\mathcal{A},\mathcal{O})$ &
    whether each assessment item corresponds to an objective &
    reverse-direction diagnostic; a high value may indicate either good coverage
    or mere objective repetition \\
    \midrule
    $\beta_{\mathrm{cong}}$ {\scriptsize(objective\_cognitive\_congruence)} &
    proportion of objectives whose best-matching assessment has the same Bloom
    level & objective--assessment cognitive-level agreement &
    assessment demand should match objective demand
    ~\cite{anderson2001taxonomy} \\
    $\beta_{\mathrm{port}}$ {\scriptsize(porter\_mean)} &
    mean Porter index over topic$\times$Bloom distributions &
    agreement between the two content--demand distributions &
    Porter alignment index~\cite{porter2002measuring,fulmer2011estimating} \\
    $\beta_{\mathrm{webb}}$ {\scriptsize(webb\_bloom\_consistency)} &
    Webb cognitive-demand consistency criterion &
    demand consistency from the item perspective &
    Webb criterion~\cite{webb1997criteria} \\
    \bottomrule
  \end{tabular}
\end{table*}

Both underlying instruments are version-pinned. The embedding encoder maps each
text span to vector $e_u$ for the $\mathrm{sim}$ operator; the primary
configuration uses Nemotron-8B, and \Cref{tab:rq2-arms} specifies alternative
encoder arms. The Bloom classifier assigns an Anderson--Krathwohl level to each
objective and assessment item for the three Family B signals; the primary
configuration is a BERT checkpoint fine-tuned on an expert-labeled learning-
objective corpus~\cite{li2022bloomclassification}, with a verb lexicon as the
comparison arm (\Cref{tab:rq2-arms}).

Finally, we report the seven signals as a pre-specified flat panel with no primary
endpoint. There are two reasons: constructive alignment has multiple facets, and
any composite would impose weights without a theoretical basis. The two
mechanically independent instrument families also enable triangulation: conclusions
are stronger only when different mechanisms agree. Signal names describe the
computed operation; ``alignment'' names only the protocol and the predicate in
\Cref{eq:aligned}, and no signal is called ``precision.'' We therefore claim only
measurable signals of constructive alignment. The protocol does not itself validate
the construct; that bridge requires a human-expert evaluation layer.
Scoring runs offline on saved outputs; instruments are configured by preset
(\Cref{tab:rq2-arms}): the primary
configuration pairs the Nemotron encoder with the fine-tuned BERT checkpoint,
and four sensitivity arms form a star around it---each arm deviates on exactly
one axis, not a grid. The BERT checkpoint always returns an argmax whereas the
lexicon can abstain, so Family-B denominators may differ between the two Bloom
arms. The panel interpretation rule is pre-specified: a size shows an advantage
when at least four of the six directional signals are significantly positive
after within-signal Holm correction and no signal is significantly negative;
the reversed pattern indicates disadvantage, and all other patterns are mixed.
The seventh, non-directional signal is still reported. Operationally, scoring
uses no LLM and has no stochastic component; results are reproducible given
pinned encoder versions and the embedding cache. Cost telemetry reads only
\texttt{metadata.token\_usage}.

\begin{table}[t]
  \caption{Instrument configurations: the primary configuration and four
  sensitivity arms (a star---each arm deviates from the primary on exactly one
  axis).}
  \label{tab:rq2-arms}
  \small
  \begin{tabular}{@{}lllp{0.2\columnwidth}@{}}
    \toprule
    Arm & Embedding encoder & Bloom classifier & Tested axis \\
    \midrule
    Primary & Nemotron-8B~\cite{babakhin2025nemotron} &
    fine-tuned BERT~\cite{li2022bloomclassification} & --- \\
    BGE-M3 & BGE-M3~\cite{chen2024m3embedding} & fine-tuned BERT & encoder lineage \\
    Granite & Granite-311M~\cite{awasthy2026granite} & fine-tuned BERT & encoder lineage \\
    TF--IDF & TF--IDF char $n$-gram & fine-tuned BERT & non-semantic floor \\
    Lexicon & Nemotron-8B & verb lexicon~\cite{anderson2001taxonomy} &
    Bloom classifier \\
    \bottomrule
  \end{tabular}
\end{table}

\subsection{Statistical analysis}
\label{sec:stats}

The unit of inference is a scenario's mean over three runs. For the level part
of Objective~2, paired Wilcoxon signed-rank tests compare A0 with each baseline at
every size and on all three judge signals; Objective~3 applies the same test to each
panel signal. When a paired sample has $n\leq25$, we use the exact permutation
null. We report paired mean differences, rank-biserial effects, and bootstrap
95\% confidence intervals from 10{,}000 scenario resamples with seed 42.

Holm correction covers the six proposed-versus-baseline comparisons per size,
separately for each signal. Objective~3 additionally reports panel-wide Holm correction
over seven signals and six comparisons as multiplicity sensitivity; this
correction creates no hierarchy among signals. For the scale part of Objective~2, let
$\Delta_s$ be the paired A0-minus-baseline gap at size $s$. The pre-specified
interaction is $\Delta_{\mathrm{largest}}-\Delta_{\mathrm{smallest}}$, tested
with paired Wilcoxon and a bootstrap interval; a negative value means a larger
A0 advantage at the small end. Objective~3 uses the same interaction, with
$\sigma_{\mathrm{meas}}$ leading the scale narrative and the other signals
reported alongside it.

For missing outputs, the primary layer assigns \texttt{failure=0}: an agent
that produces no output has quality zero, not missing data. Complete-case
inference is only a sensitivity analysis and never unconditional, because it
conditions on a post-treatment variable: it removes precisely the most severe
failures, especially at smaller sizes. Objective~2 additionally reports
\texttt{failure=min}; Objective~3 does not use this layer because its signals already
have floor zero. Output-failure rates are compared separately with Fisher's
exact test under Holm correction.

\subsection{Reproducibility}
\label{sec:reproducibility}

All measured values are generated from saved run, judge, alignment, and
ablation artifacts and enter the paper only through generated files; no number
is hand-typed. Every table cell pairs its mean with the between-run SD of the
same signal. Manifests pin runtime and model settings; revision-aware embedding
caches support exact re-analysis from saved vectors. Results produced before
the study pivot are excluded.
